\section{字符串高级算法：参数、调用与改造}

本节紧跟前面的基础模板。先看每个接口传什么，再看哪些不变量不能动，最后看如何替换维护信息或转移。

\subsection{Booth 最小表示法}

循环字符串的最小表示不是把字符串复制后排序所有起点。Booth 双指针在 \passthrough{\lstinline!s+s!} 上比较两个候选起点，淘汰整段，复杂度 \(O(n)\)。

\begin{lstlisting}[language={C++}]
// 接口：booth(s)：返回字符串字典序最小循环表示的 0-based 起点；返回类型 int。
// 参数：输入字符串 s。
int booth(const string& s) {
    // t=s+s，把所有循环移位统一成 t 中长度 n 的子串。
    // 变量：t=用于寻找最小循环表示的双倍字符串 t=s+s。
    string t = s + s;
    // 变量：n=容器 s 的元素数量 n；i=最小循环表示算法当前第一个候选起点 i；j=最小循环表示算法当前第二个候选起点 j；k=比较两个循环表示时已匹配的偏移长度 k。
    int n = s.size(), i = 0, j = 1, k = 0;
    while (i < n && j < n && k < n) {
        // k 表示两个候选起点已经匹配的长度。
        if (t[i+k] == t[j+k]) { ++k; continue; }

        // 第一个不同字符决定哪个起点更大；较大的候选可以淘汰一整段。
        if (t[i+k] > t[j+k]) i = i + k + 1;
        else j = j + k + 1;

        // 两个指针不能相等，否则候选数会少一个。
        if (i == j) ++j;
        k = 0;
    }

    // 最小表示的起点一定在 [0,n-1] 中。
    return min(i, j);
}
\end{lstlisting}

\subsubsection{测试样例：最小表示的重复字符、周期和首尾交叉}

下面每组输入一个字符串，输出它的字典序最小循环移位。

\begin{lstlisting}
输入
4
a
aaaa
baca
caba

输出
a
aaaa
abac
abac
\end{lstlisting}

前两组检查 \passthrough{\lstinline!i==j!} 和重复字符；后两组的最优起点不在原串开头，能抓住循环串复制后排序时的下标错误。

例题：\href{https://www.luogu.com.cn/problem/P1368}{P1368 工艺}。把输入串交给 \passthrough{\lstinline!booth!}，从返回下标开始输出 \passthrough{\lstinline!n!} 个字符。题目求的是循环移位，不是反转后的最小串；若题面允许反转，要额外对反串再调用一次。


\subsection{回文自动机（Palindromic Automaton）}

PAM 的每个节点代表一个不同的回文串；\passthrough{\lstinline!fail!} 指向去掉两端后最长的回文后缀。插入字符时沿 \passthrough{\lstinline!fail!} 找到可以向两端扩展的节点。节点数最多 \passthrough{\lstinline!n+2!}，所以它适合``所有回文子串''的统计。

\begin{lstlisting}[language={C++}]
struct PAM {
    static const int SIG = 26; // 字符集大小；本模板要求输入字符在 'a'..'z'。
    int next[MAXN][SIG]; // next[u][c]：回文串 u 两端加字符 c 后的节点。
    int fail[MAXN];       // fail[u]：u 的最长真回文后缀。
    int len[MAXN];        // len[u]：节点 u 代表的回文串长度。
    int s[MAXN];          // 已插入的字符，get_fail 要访问历史字符。
    long long occ[MAXN];  // 先统计“作为最长回文后缀”的次数，再沿 fail 汇总。
    int last, tot, n;     // last=当前前缀的最长回文后缀，tot=节点总数，n=串长。

    // 接口：init()：清空成员并建立本自动机所需的初始根节点；完成初始化；多测时每组重新调用。
    void init() {
        // 0 号根表示长度 -1 的虚拟回文，1 号根表示长度 0 的空串。
        memset(next, 0, sizeof(next));
        memset(occ, 0, sizeof(occ));
        n = 0; tot = 1; last = 1;
        len[0] = -1; fail[0] = 0;
        len[1] = 0; fail[1] = 0;
    }
    // 接口：get_fail(x)：沿 fail 链找到能被当前字符继续扩展的节点；返回类型 int。
    // 参数：开始沿 fail 链检查的自动机状态编号 x。
    int get_fail(int x) {
        // 从当前回文后缀沿 fail 向上跳，直到它能和新字符首尾匹配。
        while (n - 1 - len[x] < 0 || s[n - 1 - len[x]] != s[n])
            x = fail[x];
        return x;
    }
    // 接口：add(ch)：向回文自动机追加字符，更新 last 并返回新的最长回文后缀状态编号；会修改当前结构；多测时重新构造或显式清空成员。
    // 参数：待加入字符 ch。
    int add(char ch) {
        // 新字符暂存在 s[n]，此时 n 还是旧长度。
        s[n] = ch - 'a';
        // 变量：cur=沿 fail 链找到的、可被新字符扩展的最长回文后缀状态 cur；c=本次加入回文自动机的字符编号 c。
        int cur = get_fail(last), c = s[n];
        if (!next[cur][c]) {
            // 当前回文后缀加上两端字符，产生一个此前没有出现过的新回文。
            // 变量：now=本次新建的回文自动机状态编号 now。
            int now = ++tot;
            len[now] = len[cur] + 2;

            // 长度 1 的回文 fail 到空串根；其它节点继续从 fail[cur] 找。
            fail[now] = (len[now] == 1) ? 1 : next[get_fail(fail[cur])][c];
            next[cur][c] = now;
        }

        // 新前缀的最长回文后缀就是这条转移到达的节点。
        last = next[cur][c]; ++occ[last]; ++n;
        return last;
    }
    // 接口：count_occurrence()：按回文长度从大到小把 occ 沿 fail 链累加，原地得到每个回文状态的出现次数；无返回值。
    void count_occurrence() {
        // fail 一定从长串指向短串，所以必须按 len 从大到小传递。
        // 节点编号不保证按长度排序，不能盲写 for(i=tot;i>=1;--i)。
        // 变量：ord=按回文长度从大到小排列、用于沿 fail 链汇总 occ 的回文自动机状态编号 ord。
        vector<int> ord(max(0, tot - 1));
        iota(ord.begin(), ord.end(), 2); // 0/1 是两个虚拟根，不参与统计。
        // 接口：比较器 lambda(a, b)：按回文长度从大到小排列自动机状态，供出现次数沿 fail 链汇总。
        // 参数：第一个回文状态 a；第二个回文状态 b。
        sort(ord.begin(), ord.end(), [&](int a, int b) {
            return len[a] > len[b];
        });
        // 变量：u=按指定顺序枚举的自动机状态编号 u。
        for (int u : ord) occ[fail[u]] += occ[u];
    }
};
\end{lstlisting}

\subsubsection{测试样例：回文自动机的根、重复回文和多测清空}

下面每组输入一个小写字符串，输出不同的非空回文子串数量，也就是 \passthrough{\lstinline!tot-1!}。

\begin{lstlisting}
输入
4
a
aaaa
abacaba
ababa

输出
1
4
7
5
\end{lstlisting}

\passthrough{\lstinline!aaaa!} 的不同回文串是 \passthrough{\lstinline!a,aa,aaa,aaaa!}；\passthrough{\lstinline!abacaba!} 的七个不同回文串可用来检查 \passthrough{\lstinline!fail!} 链。连续多组调用前必须执行 \passthrough{\lstinline!init()!}，否则第二组的节点和出现次数会污染第三组。

例题：\href{https://www.luogu.com.cn/problem/P5496}{P5496 【模板】回文自动机（PAM）}。题目要求每个位置结尾的回文串数量；插入第 \passthrough{\lstinline!i!} 个字符后，最长回文后缀是 \passthrough{\lstinline!last!}，所有回文后缀沿 \passthrough{\lstinline!fail!} 链连接，因此可以用 \passthrough{\lstinline!dp[i] = dp[i-len[last]] + 1!} 的题目特定统计，或沿 fail 链统计。若只问不同回文子串数量，答案是 \passthrough{\lstinline!tot-1!}；若问出现次数，要先 \passthrough{\lstinline!count\_occurrence()!}。


\subsection{后缀数组、最长公共前缀、后缀自动机与 Aho--Corasick 失配树}

后缀数组适合``后缀排序、子串字典序、多个后缀的最长公共前缀''；LCP 用 Kasai 求相邻后缀公共前缀，再用 RMQ 回答任意两个后缀的 LCP。题目出现``第 k 个不同子串''时，SAM 往往更直接：每个状态记录 \passthrough{\lstinline!len!}、\passthrough{\lstinline!link!} 和转移，按 DAG DP 统计从状态出发的不同串数。

\begin{lstlisting}[language={C++}]
struct SAM {
    int next[MAXN][26]; // 状态转移；字符集大时改成 map/unordered_map。
    int link[MAXN];     // 后缀链接：指向最长的更短 endpos 等价类。
    int len[MAXN];      // 该状态能表示的最长子串长度。
    int last = 1;       // 整个当前前缀对应的状态。
    int tot = 1;        // 1 号状态是初始状态；0 号转移为空。

    // 接口：extend(c)：向当前字符串自动机加入一个字符并更新状态；会修改当前结构；多测时重新构造或显式清空成员。
    // 参数：待加入自动机的 0-based 字母表编号 c。
    void extend(int c) {
        // 新状态 cur 代表“以新字符结尾”的所有后缀。
        // 变量：cur=加入新字符后创建的 SAM 状态编号 cur。
        int cur = ++tot; len[cur] = len[last] + 1;
        // 变量：p=SAM 扩展前原末状态的游走指针 p。
        int p = last;

        // 沿 suffix link 补上以前不存在的 c 转移。
        while (p && !next[p][c]) next[p][c] = cur, p = link[p];
        if (!p) link[cur] = 1;
        else {
            // 变量：q=SAM 状态 p 经当前字符转移到的既有状态 q。
            int q = next[p][c];
            if (len[p] + 1 == len[q]) link[cur] = q;
            else {
                // q 的长度太长，不能直接作为 cur 的 link。
                // clone 复制 q 的转移，但把最长长度截为 len[p]+1。
                // 变量：clone=后缀自动机为拆分转移而复制的克隆状态 clone。
                int clone = ++tot;
                memcpy(next[clone], next[q], sizeof next[q]);
                len[clone] = len[p] + 1; link[clone] = link[q];

                // 把原来指向 q 的、位于同一等价区间的转移改指 clone。
                while (p && next[p][c] == q)
                    next[p][c] = clone, p = link[p];
                link[q] = link[cur] = clone;
            }
        }
        last = cur;
    }
};
\end{lstlisting}

\subsubsection{调用测试：后缀自动机的空串、重复字符和克隆状态}

当前代码片段只负责 \passthrough{\lstinline!extend!}，所以先测试最基本的不变量：初始状态编号为 \passthrough{\lstinline!1!}，字符串全相同时每个字符新增一个状态；如果要测试出现次数，还必须把 \passthrough{\lstinline!occ!} 统计和链接树排序一并接上。

\begin{lstlisting}
调用输入
s = ""
s = "a"
s = "aaaa"
s = "ab"

应得到的状态总数 tot
1
2
5
3
\end{lstlisting}

第一组检查 \passthrough{\lstinline!init!}；第二、三组检查 \passthrough{\lstinline!last!} 和长度；第四组检查不同字符转移。真正测试克隆分支时，应使用一个会产生重复后缀等价类的字符串，并检查所有被旧状态指向的转移是否一起改到 clone。

例题：\href{https://www.luogu.com.cn/problem/P3804}{P3804 【模板】后缀自动机（SAM）}问出现次数最多的子串：按 \passthrough{\lstinline!len!} 从大到小把 \passthrough{\lstinline!occ[u]!} 加到 \passthrough{\lstinline!occ[link[u]]!}，答案是 \passthrough{\lstinline!max(len[u]*occ[u])!}。若是第 \passthrough{\lstinline!k!} 个子串，先在状态 DAG 上做 \passthrough{\lstinline!dp[u]=sum(1+dp[v])!}，再按字符顺序跳转并扣减 \passthrough{\lstinline!1+dp[v]!}。

AC 自动机把多个模式串建 Trie，再按 BFS 建 \passthrough{\lstinline!fail!}；把 fail 指针反向成树，就能把``文本中经过节点''转成 fail 子树求和。例题可用\href{https://www.luogu.com.cn/problem/P3796}{P3796 AC 自动机加强版}：插入模式串，扫描文本，沿 fail 链累计出现次数，最后找最大次数的模式串。查询中如果有``一个模式串是另一个模式串后缀''，不要只看 Trie 节点本身，必须沿 fail 树汇总。

哈希二分 LCP 的模板：预处理双模/64 位前缀哈希，比较 \passthrough{\lstinline!hash(i,len)==hash(j,len)!}，二分最长长度。例题：\href{https://www.luogu.com.cn/problem/P3370}{P3370 【模板】字符串哈希}。如果是高风险哈希题，使用双模或 \passthrough{\lstinline!uint64\_t!} 随机基，不能把单模碰撞当成严格证明。


\subsection{后缀数组与 LCP}

后缀数组的倍增法按``前半排名、后半排名''排序。工程上可以使用 \passthrough{\lstinline!sort!} 版先保证正确；\passthrough{\lstinline!n!} 达到 \passthrough{\lstinline!2e5!} 以上再换基数排序优化。

\begin{lstlisting}[language={C++}]
// 接口：suffix_array(s)：返回 sa；sa[i] 是字典序第 i 小后缀的 0-based 起点；返回类型 vector<int>。
// 参数：待构造后缀数组的字符串 s。
vector<int> suffix_array(const string& s) {
    // 变量：n=容器 s 的元素数量 n。
    int n = s.size();
    if (n == 0) return {};
    // 变量：sa=后缀数组 sa；sa[i] 是第 i 小后缀起点；rk=rk[i] 表示后缀 s[i..] 按当前倍增长度得到的等价类排名；tmp=本轮按长度 2k 的排名二元组排序后写入的新等价类排名数组 tmp。
    vector<int> sa(n), rk(n), tmp(n);
    iota(sa.begin(), sa.end(), 0);
    // rk[i] 是长度 1 的后缀 s[i..] 的排名；不同字符最好先压成连续整数。
    // 变量：i=后缀数组初始化单字符排名的后缀起点 i。
    for (int i = 0; i < n; ++i) rk[i] = s[i];
    // 变量：k=后缀倍增中当前二元组每一段的长度 k；本轮比较长度 2k 的前缀。
    for (int k = 1;; k <<= 1) {
        // 现在按长度 2k 的二元组 (前 k 长排名, 后 k 长排名) 排序。
        // 接口：比较器 lambda(i, j)：按长度 2k 的两段排名二元组比较后缀 i 与 j。
        // 参数：第一个后缀起点 i；第二个后缀起点 j。
        sort(sa.begin(), sa.end(), [&](int i, int j) {
            if (rk[i] != rk[j]) return rk[i] < rk[j];
            // 变量：ri=后缀 i 偏移 k 后第二段的排名 ri；越界时为 -1。
            int ri = i + k < n ? rk[i+k] : -1;
            // 变量：rj=后缀 j 偏移 k 后第二段的排名 rj；越界时为 -1。
            int rj = j + k < n ? rk[j+k] : -1;
            return ri < rj;
        });
        // 根据排好序的后缀重新编号；相邻二元组相同才共享排名。
        tmp[sa[0]] = 0;
        // 变量：i=后缀数组排序后用于比较相邻后缀并重编号的排名位置 i。
        for (int i = 1; i < n; ++i) {
            // 变量：a=字典序前一个后缀的起点 a=sa[i-1]；b=字典序当前后缀的起点 b=sa[i]。
            int a = sa[i-1], b = sa[i];
            // 变量：diff=相邻后缀 a、b 的第一段长度 k 排名是否不同的标记 diff。
            bool diff = rk[a] != rk[b];
            // 变量：ra=后缀 a 偏移 k 后第二段的排名 ra；越界时为 -1。
            int ra = a+k < n ? rk[a+k] : -1;
            // 变量：rb=后缀 b 偏移 k 后第二段的排名 rb；越界时为 -1。
            int rb = b+k < n ? rk[b+k] : -1;
            tmp[b] = tmp[a] + (diff || ra != rb);
        }
        rk.swap(tmp);
        // 所有排名都不同后，sa 已经是最终后缀数组。
        if (rk[sa.back()] == n - 1) break;
    }
    return sa;
}
// 接口：kasai(s, sa)：由 s 和 sa 返回相邻后缀 LCP 数组；返回类型 vector<int>。
// 参数：后缀数组对应的原字符串 s；原字符串的后缀数组 sa。
vector<int> kasai(const string& s, const vector<int>& sa) {
    // 变量：n=容器 s 的元素数量 n；rank=每个后缀起点在后缀数组 sa 中的排名 rank；lcp=相邻后缀最长公共前缀数组 lcp。
    int n = s.size(); vector<int> rank(n), lcp(n);
    // 变量：i=后缀数组 sa 中当前排名位置 i。
    for (int i = 0; i < n; ++i) rank[sa[i]] = i;
    // 变量：i=Kasai 算法当前计算 LCP 的后缀起点 i。
    for (int i = 0, h = 0; i < n; ++i) {
        // 变量：r=后缀 i 在后缀数组 sa 中的排名 r。
        int r = rank[i];
        if (r == 0) continue;
        // 变量：j=后缀数组中排名紧邻后缀 i 之前的后缀起点 j。
        int j = sa[r-1];
        while (i+h < n && j+h < n && s[i+h] == s[j+h]) ++h;
        lcp[r] = h;
        if (h) --h;
    }
    return lcp; // lcp[r] = sa[r] 与 sa[r-1] 的 LCP
}
\end{lstlisting}

\subsubsection{测试样例：后缀排序和 LCP 的经典字符串}

\begin{lstlisting}
输入字符串
banana

应得到后缀数组
5 3 1 0 4 2

应得到 lcp 数组
0 1 3 0 0 2
\end{lstlisting}

\passthrough{\lstinline!banana!} 同时含有重复前缀和不同首字符，能检查倍增第二关键字的越界排名；\passthrough{\lstinline!lcp[2]=3!} 是 \passthrough{\lstinline!ana!} 与 \passthrough{\lstinline!anana!} 的公共前缀，能抓住 Kasai 的 \passthrough{\lstinline!h!} 是否正确复用。

例题：\href{https://www.luogu.com.cn/problem/P3809}{P3809 【模板】后缀排序}。若要任意两个后缀的 LCP，把 \passthrough{\lstinline!lcp!} 做 RMQ；若要第 \passthrough{\lstinline!k!} 个不同子串，按后缀顺序累加 \passthrough{\lstinline!n-sa[i]-lcp[i]!}，第一次超过 \passthrough{\lstinline!k!} 的后缀就是答案所在位置。


\subsection{AC 自动机与 fail 树}

\begin{lstlisting}[language={C++}]
struct AC {
    int ch[MAXN][26]{}; // build 后不存在的转移也会补成 fail 转移。
    int fail[MAXN]{};   // 当前节点的最长真后缀节点。
    int out[MAXN]{};    // 节点本身及 fail 链上的模式串计数。
    int tot = 0;          // 当前最后一个节点编号；根为 0，节点范围 [0,tot]。
    vector<int> tree[MAXN]; // fail 指针反图；用于子树汇总。
    // 接口：insert(s)：把小写模式串 s 插入 Trie，并在终止节点累计出现次数；会修改当前结构；多测时重新构造或显式清空成员。
    // 参数：输入字符串 s。
    int insert(const string& s) {
        // 变量：u=字符串自动机从根开始游走的状态编号 u。
        int u = 0;
        // 变量：cc=从模式串 s 依次读取的字符 cc。
        for (char cc : s) {
            // 变量：c=当前文本字符映射到的字母表编号 c。
            int c = cc - 'a';
            if (!ch[u][c]) ch[u][c] = ++tot;
            u = ch[u][c];
        }
        // 返回结束节点，题目要分别输出时保存这个编号。
        return u;
    }
    // 接口：build()：用 BFS 建立 AC 自动机 fail 指针、补全转移并汇总输出计数；完成初始化；多测时每组重新调用。
    void build() {
        // 变量：q=AC 自动机建立 fail 指针时按层遍历状态的 BFS 队列 q。
        queue<int> q;
        // 根的 fail 是自己；根的直接儿子 fail 都是根。
        // 变量：c=当前枚举的 0-based 字母表编号 c。
        for (int c = 0; c < 26; ++c)
            if (ch[0][c]) q.push(ch[0][c]);
        while (!q.empty()) {
            // 变量：u=AC 自动机构建过程中当前出队并补全转移的状态编号 u。
            int u = q.front(); q.pop();
            // 变量：c=当前枚举的 0-based 字母表编号 c。
            for (int c = 0; c < 26; ++c) {
                if (ch[u][c]) {
                    // 子节点的 fail：沿 u 的 fail 走同一个字符。
                    fail[ch[u][c]] = ch[fail[u]][c];
                    q.push(ch[u][c]);
                } else {
                    // 补全自动机转移，扫描文本时就不需要 while 回退。
                    ch[u][c] = ch[fail[u]][c];
                }
            }
        }
        // fail[u] 是 u 的父亲，之后可以把文本经过次数沿 fail 树向上累加。
        // 变量：u=当前加入 fail 树的 AC 自动机状态编号 u。
        for (int u = 1; u <= tot; ++u) tree[fail[u]].push_back(u);
    }
    // 接口：scan(text)：扫描文本 text，并把每次到达自动机状态的次数原地累加到 occ；无返回值。
    // 参数：主串 text。
    void scan(const string& text) {
        // 变量：u=字符串自动机从根开始游走的状态编号 u。
        int u = 0;
        // 变量：cc=从主串 text 依次读取并送入 AC 自动机的字符 cc。
        for (char cc : text) {
            // 每读一个字符，u 就是当前前缀的最长 Trie 后缀状态。
            u = ch[u][cc - 'a'];
            ++out[u];
        }
    }
    // 接口：dfs_fail(u)：沿 fail 树汇总模式串出现次数；不返回值。
    // 参数：顶点编号 u。
    void dfs_fail(int u) {
        // 变量：v=从 tree[u] 枚举的相邻顶点 v。
        for (int v : tree[u]) {
            dfs_fail(v);
            // v 的所有出现也意味着 fail[v] 对应的模式串出现。
            out[u] += out[v];
        }
    }
};
\end{lstlisting}

\subsubsection{测试样例：AC 自动机的后缀模式和重复出现}

模式串依次为 \passthrough{\lstinline!a!}、\passthrough{\lstinline!ab!}、\passthrough{\lstinline!bab!}，文本为 \passthrough{\lstinline!abab!}。扫描后沿 fail 树汇总，三个模式的出现次数应为：

\begin{lstlisting}
a   -> 2
ab  -> 2
bab -> 1
\end{lstlisting}

\passthrough{\lstinline!ab!} 是 \passthrough{\lstinline!a!} 的延长，\passthrough{\lstinline!bab!} 又是文本中较长的后缀模式；如果只读取当前 Trie 节点的 \passthrough{\lstinline!out!} 而不沿 fail 树汇总，\passthrough{\lstinline!a!} 的次数通常会少算。

例题：\href{https://www.luogu.com.cn/problem/P3808}{P3808 【模板】AC 自动机（简单版）}。每个模式串的结束节点记录 \passthrough{\lstinline!id!}；扫描文本后，调用 \passthrough{\lstinline!dfs\_fail(0)!}，该节点的 \passthrough{\lstinline!out!} 就是模式串作为后缀出现的总次数。简单版只需问总出现次数，若要每个模式串分别输出，就保存插入时的节点编号。
